\scanpage{69}
\mspage{30}
\uline{propriété universelle} évidente pour les morphismes
de la catégorie topologique $\uline{\Sigma}$ dans un ens. \inserted{pré}ordonné.
\unsure{\S} Supposons maintenant que $\uline{\Sigma}$ satisfasse les
conditions
\inserted{\unsure{c}}
\deleted{a) toute flèche inversible de $\uline{\Sigma}$ est
une identité}
\deleted{a) L'application $\delta:\Sigma_0\longrightarrow\Sigma_1$ est
un plongement ouvert et fermé. C'est le cas
si p.ex. $\uline{b}:\Sigma_1\longrightarrow\Sigma_0$ est une immersion
propre\ldots.}

a) $\Sigma_{(1)}\xrightarrow{\uline{b}}\Sigma_{(0)}$ est une immersion
propre ($\Longrightarrow\ \delta:\Sigma_{(0)}\longrightarrow\Sigma_{(1)}$ est une imm.
ouverte et fermée, en sorte que \deleted{$\Sigma_{(1)}=\delta(\Sigma_{(0)})\amalg$}
\[
\Sigma_{(1)}\simeq\delta(\Sigma_{(0)})\amalg\Sigma^{*}(1)
\]
)

b) $\Sigma^{*}_{(1)}\xrightarrow{\uline{b}|\Sigma^{*}(1)}\Sigma_{(0)}$ est une immersion à image
localt \uline{rare} i.e. à intérieur vide\ldots

c) \deleted{\unsure{---}} Toute composante connexe de $\Sigma_{(0)}$ est de
dim finie.
Alors on trouve que $\Sigma_{ii}^{*}\overset{\mathrm{déf}}{=}\Sigma_{ii}\cap\Sigma^{*}(1)=\varnothing$,
par des raisons de dimension (car $\Sigma_{ii}^{*}$ est
fibré sur \deleted{sur} $\Sigma_i$, donc s'il est non vide, de
dim $\geq\dim\Sigma_i$; mais comme $\Sigma_{ii}^{*}\longrightarrow\Sigma_i^{*}$ est
une imm. v.v., on doit avoir aussi $\dim\Sigma_{ii}^{*}<$
$\dim\Sigma_i^{*}$, donc $\unsure{\Sigma_{ii}^{<}|\Sigma_{ii}^{*}}$ est vide), i.e.
\[
\boxed{\Sigma_i\xrightarrow{\sim}\Sigma_{ii}}:
\]
toute flèche \deleted{dans} entre deux
\deleted{\unsure{objets de}} éléments de $\Sigma_{(0)}$ qui sont dans une
même composante connexe \deleted{sont} est une identité.
\mspage{68}
